% Method section, v1 — Continuous GRPO + Adaptive-O2O scheduler.
% Every implementation claim cites file:line in the repo. Where the prompt's
% formulation diverges from the code, the code wins and a TODO comment flags
% the deviation.

\section{Method}
\label{sec:method}

\subsection{Problem setup}

We allocate capital across $N{=}8$ ETFs (SPY, EEM, TLT, HYG, DBC, GLD, UUP,
SHY) at daily resolution. The agent emits a portfolio weight vector
$w_t \in \Delta^{N-1}$ (the $(N{-}1)$-simplex). One env step yields the
post-cost log return
\begin{equation}
\label{eq:reward}
r_t \;=\; \log\!\big(1 + \max(w_t^{\top} R_{t+1} - \lambda\, \|w_t - w_{t-1}\|_1,\, -0.9999)\big),
\end{equation}
where $R_{t+1}$ is the realized one-step return vector and $\lambda$ the
linear transaction-cost coefficient. This is implemented in
\texttt{core/envs/portfolio\_env.py:155} and re-derived bit-for-bit in
\texttt{peek\_reward\_batch} (line~249), which is the operation that makes
group sampling cheap (Sec.~\ref{sec:method:grpo}).

The state $s_t$ is a 216-dimensional vector of causal market features --- past
log-returns, rolling volatility, multi-window moving averages, and momentum
signals --- with no portfolio-state coordinates. The lack of a $w_{t-1}$
component in the observation is structural, not stylistic: it is what makes
the exogeneity assumption hold (Sec.~\ref{sec:method:grpo:exogeneity}). The
optional flag \verb|include_prev_weights=True|
(\texttt{core/envs/portfolio\_env.py:72}) injects the previous weight vector
into the observation and breaks exogeneity; we leave it disabled, and
\texttt{tests/test\_exogeneity.py:104} guards the constraint with a unit test.

\subsection{Continuous GRPO}
\label{sec:method:grpo}

Group Relative Policy Optimization \citep{shao2024deepseekmath} was
introduced for discrete language-model action spaces: at each prompt
sequence, $G$ candidate completions are sampled, and a per-prompt
advantage is computed by within-group normalization, sidestepping the need
for a learned value baseline. We carry the same group-relative idea into
the continuous portfolio simplex, leveraging an environment-level property
that LM training does not enjoy: the next state is independent of the
agent's action.

\subsubsection{Policy parameterization: Dirichlet head on the simplex}

The actor is a 2-layer MLP with $\tanh$ activations whose output head emits
$N$ concentration parameters
($\texttt{core/networks/policies.py:401--429}$):
\begin{equation}
\alpha(s_t) \;=\; \softplus(W_\alpha h(s_t) + b_\alpha) + \mathbf{1},
\quad
\pi_\theta(\,\cdot \mid s_t) \;=\; \mathrm{Dirichlet}(\alpha(s_t)).
\end{equation}
The $+1$ shift forces $\alpha_i > 1$, making the Dirichlet density unimodal
and avoiding the simplex-corner singularities that vanilla softmax-of-Gaussian
heads suffer near $w_i \to 0$. Sampling, log-probability, and entropy are
PyTorch's closed-form Dirichlet calls --- no projection step is needed
because samples are already in the open simplex (modulo a numerical
$10^{-6}$ clamp at \texttt{core/networks/policies.py:452}).

\subsubsection{Group sampling under exogenous state transitions}
\label{sec:method:grpo:exogeneity}

At rollout time, at each visited state $s_t$ we draw
$\{a_t^{(i)}\}_{i=1}^G \sim \pi_{\theta_{\text{old}}}(\cdot \mid s_t)$ from
the current actor, then compute their per-candidate rewards
$\{r_t^{(i)}\}_{i=1}^G$ from the same realized return $R_{t+1}$
(\texttt{online\_rl/agents/grpo.py:271--284}). This is the \emph{group}
in GRPO: $G$ siblings sharing $s_t$ and $R_{t+1}$ but differing in $a_t^{(i)}$.

The soundness of using these per-candidate rewards as relative advantages
hinges on a property of the environment, not of the policy:
\begin{equation}
\label{eq:exogeneity}
P(s_{t+1} \mid s_t, a_t) \;=\; P(s_{t+1} \mid s_t) \qquad \text{(exogeneity)}.
\end{equation}
A retail allocation does not move broad-index ETF prices on a one-day
horizon, so the next observation depends only on the realized market path,
not the agent's weight choice. The default-obs PortfolioEnv satisfies
Eq.~\ref{eq:exogeneity} exactly --- no market-impact model, no
weight-aware feature, no replay-style state mutation. We verify this with
two adversarial unit tests
(\texttt{tests/test\_exogeneity.py:52, 91}) that step two seed-aligned envs
with maximally different one-hot allocations and assert bit-identical
$s_{t+1}$, and a third test
(\texttt{tests/test\_exogeneity.py:104}) that documents the configuration
flag which would break the property.

Under Eq.~\ref{eq:exogeneity}, the standard advantage
$A^{(i)} = Q(s_t, a_t^{(i)}) - V(s_t)$ admits a critic-free closed form
\emph{within a group}: the bootstrapped continuation value
$\gamma\, \mathbb{E}[V(s_{t+1}) \mid s_t, a_t^{(i)}] = \gamma\, \mathbb{E}[V(s_{t+1}) \mid s_t]$
is independent of $i$ and cancels under within-group centering. What
remains is just the immediate reward difference, which we score directly
through \texttt{env.peek\_reward\_batch} --- a method that recomputes the
step reward formula (Eq.~\ref{eq:reward}) for $G$ candidates against the
same $R_{t+1}$ pointer without mutating env state
(\texttt{core/envs/portfolio\_env.py:195--250}).

\subsubsection{Group-relative advantage}

The advantage of candidate $i$ at state $s_t$ is the within-group z-score:
\begin{equation}
\label{eq:grpo-adv}
A_t^{(i)} \;=\; \frac{r_t^{(i)} - \bar{r}_t}{\sigma_{r_t} + \varepsilon},
\quad
\bar{r}_t = \tfrac{1}{G}\sum_j r_t^{(j)},
\quad
\sigma_{r_t}^2 = \tfrac{1}{G}\sum_j (r_t^{(j)} - \bar{r}_t)^2,
\end{equation}
with $\varepsilon = 10^{-8}$. This is the default
\verb|advantage_norm="mean_std"| mode in
\texttt{online\_rl/agents/grpo\_advantages.py:66--69}. The same module also
implements three ablation variants (\verb|"raw"|, \verb|"mean_only"|,
\verb|"rank"|; lines~60--72, 75--95) that we use for the advantage-norm
ablation in Sec.~\ref{sec:experiments:phase2d}.

A subtle but load-bearing detail: when all $G$ candidates score identical
rewards (e.g., all collapse to equal-weight on a flat-return day), the
$\sigma_r{+}\varepsilon$ denominator keeps the advantage finite and the
gradient zero. This prevents the constant-reward pathology that a naive
$(r - \bar r)/\sigma_r$ would produce.

\subsubsection{Loss}

We fit $\pi_\theta$ with a PPO-style clipped surrogate plus a KL anchor to
a frozen reference $\pi_{\text{ref}}$ snapshotted at training start
(\texttt{online\_rl/agents/grpo.py:227--231}):
\begin{equation}
\label{eq:grpo-loss}
\mathcal{L}(\theta) \;=\;
\underbrace{- \mathbb{E}_{s, i}\!\left[\min\!\left(
\rho^{(i)} A^{(i)},\;
\mathrm{clip}(\rho^{(i)}, 1{-}\epsilon, 1{+}\epsilon)\, A^{(i)}
\right)\right]}_{\text{surrogate}}
\;+\; \beta_{\mathrm{KL}}\, \mathbb{E}_s\!\left[D_{\mathrm{KL}}(\pi_\theta \,\|\, \pi_{\text{ref}})\right]
\;-\; c_H\, \mathbb{E}_s[H(\pi_\theta)],
\end{equation}
where $\rho^{(i)} = \pi_\theta(a^{(i)} \mid s) / \pi_{\theta_{\text{old}}}(a^{(i)} \mid s)$
is the importance ratio against the actor that produced the samples. The
log-ratio is hard-clamped to $[\log 10^{-3}, \log 10^{3}]$
(\texttt{online\_rl/agents/grpo.py:79--80, 145--148}) before exponentiation;
this is belt-and-suspenders against the first-epoch divergence that PPO
clipping alone admits. Defaults: $G{=}16$, $\epsilon{=}0.2$,
$\beta_{\mathrm{KL}}{=}0.01$, $c_H{=}0$, 4 epochs of size-256 minibatches per
update phase (\texttt{online\_rl/agents/grpo.py:34--71}).

The KL anchor term is the closed-form Dirichlet$\to$Dirichlet KL
(\texttt{online\_rl/agents/grpo.py:164--166}); no Monte-Carlo estimate. The
reference actor's parameters are frozen for the entire run and never
re-snapshotted --- a stronger anchor than RLHF-style periodic refreshes,
chosen so that the policy cannot drift unboundedly under the small-batch
financial signal-to-noise ratio.

% TODO(post-2B): the persona's GRPO framing leaves the importance ratio out
% of the advantage equation. The code's surrogate includes both an
% importance-ratio (rho) and a clipped-surrogate term --- closer to PPO
% than to vanilla GRPO. We describe what the code does; if the report needs
% the simpler form, drop the clipping by setting cfg.clip_eps very large
% (the hard-ratio clamp would still prevent NaN propagation).

\subsubsection{Why this is sound only here}

Continuous GRPO trades a learned critic for an environmental property. In
domains where actions \emph{do} affect $s_{t+1}$ (e.g., locomotion, robotics,
market-impact-aware finance), the bootstrap term does not cancel and group
advantages are biased. The unit tests in \texttt{tests/test\_exogeneity.py}
are not stylistic --- they are an algorithmic precondition. We surface this
explicitly in the limitations (Sec.~\ref{sec:limitations}) and recommend
that any future port of this method include an analogous test against the
host env.

\subsection{Adaptive-O2O scheduler}
\label{sec:method:o2o}

The offline-to-online (O2O) pipeline pre-trains a Geodesic-CQL critic on
2008--2020 transitions (the \emph{offline} phase) and fine-tunes a
SAC-Dirichlet actor on the 2022Q1--2026Q1 test window (the \emph{online}
phase), mixing offline and online batches in each gradient step
(\texttt{hybrid\_rl/agents/o2o\_agent.py:204--263}). The adaptive
contribution is a per-step schedule on the conservatism weight that
attenuates pessimism when the online state distribution looks like
something the offline buffer covered.

\paragraph{Regime encoder and KL.} A 1-layer GRU over the last 20-day
observation window produces a regime embedding $h_t \in \mathbb{R}^{64}$
(\texttt{hybrid\_rl/configs/o2o\_config.py:7}). At training start we record
the offline-batch encoder outputs as the reference distribution; during
fine-tuning we maintain a 5-step rolling buffer of online encoder outputs
(\texttt{hybrid\_rl/agents/o2o\_agent.py:225--241}). The regime KL is a
diagonal-Gaussian closed form fit per coordinate
(\texttt{hybrid\_rl/agents/o2o\_agent.py:23--32}):
\begin{equation}
\widehat{\mathrm{KL}}(h_{\text{off}} \,\|\, h_{\text{on}}) \;=\;
\sum_{i=1}^{64} \log\!\frac{\sigma^{\text{on}}_i}{\sigma^{\text{off}}_i}
+ \frac{(\sigma^{\text{off}}_i)^2 + (\mu^{\text{off}}_i - \mu^{\text{on}}_i)^2}{2 (\sigma^{\text{on}}_i)^2}
- \tfrac{1}{2}.
\end{equation}

\paragraph{Conservatism weight.} The CQL penalty multiplier passed into
\texttt{sac\_agent.update\_critic(\dots, cql\_weight=$\beta_t$)} is
\begin{equation}
\label{eq:cql-schedule}
\beta_t \;=\; \alpha_{\mathrm{CQL}} \cdot \sigma\!\big(\eta \cdot \widehat{\mathrm{KL}}(h_{\text{off}} \,\|\, h_{\text{on}})\big),
\end{equation}
with $\alpha_{\mathrm{CQL}} = 5.0$ and $\eta = 0.5$ as defaults
(\texttt{hybrid\_rl/configs/o2o\_config.py:22--24};
implementation \texttt{hybrid\_rl/agents/o2o\_agent.py:166--200}). Under
distribution shift the sigmoid saturates near $1$, restoring full offline
conservatism; under distribution match it relaxes toward zero, freeing the
online updates to chase the in-domain reward signal. The naive-fine-tune
condition in our four-way comparison (Sec.~\ref{sec:experiments:phase2c})
sets $\beta_t \equiv \alpha_{\mathrm{CQL}}$ throughout, exposed as
\texttt{--adaptive\_conservatism=false} in the runner.

\paragraph{The bug fix that made the comparison meaningful.} The repo
shipped with \texttt{run\_o2o.py} containing an inline online-phase loop
that called \texttt{sac\_agent.update\_critic(mixed)} \emph{without}
forwarding any $\beta_t$, defaulting the CQL weight to zero and silently
bypassing the entire scheduler we describe above. The full
\texttt{O2OAgent.finetune\_online} method --- which actually plumbs
$\beta_t$ through the critic update --- existed but was never invoked from
the CLI. We replaced the inline loop with a call to
\texttt{agent.finetune\_online(n\_online)}
(\texttt{scripts/run\_o2o.py:309--312}); without this, every Phase 2C
``adaptive O2O'' run would have been numerically indistinguishable from
the naive condition, regardless of how the CLI flag was set. We document
this fix and three other pre-existing bugs in
\texttt{writeup/draft\_implementation\_notes.md}.

\paragraph{Initialization and env handoff.} \texttt{transfer\_to\_online}
copies the actor / critic / regime-encoder / temperature parameters from
the Geodesic-CQL agent to the SAC-Dirichlet agent, then re-resets the
shared training env so the SAC agent's cached observation does not point
at a terminal state left over from offline buffer collection
(\texttt{hybrid\_rl/agents/o2o\_agent.py:142--164}; the env reset on
lines~155--164 is also a bug fix from a pre-existing crash documented in
the implementation notes).

% TODO(post-2C): if the cql\_weight\_traj.npy series shows $\beta_t$ saturating
% at one extreme for the entire test window, we will need to either widen
% the regime-encoder context or revisit $\eta$. The naive-vs-adaptive
% comparison loses signal if the schedule is constant in practice.
